\scanpage{647}
\mspage{35}

Je ne reviens pas ici sur la théorie de K. S. Brown
des catégories à cofibrations. C'est à ma connaissance
l'ensemble d'axiomes le plus économique
connu, pour assurer d'une part une bonne
prise sur $\operatorname{Ho}_W$ (par une notion convenable
\unsure{d'|de'}homotopie d'applications, et un calcul \unsure{de|des}
fractions dans $\pi M_c$), et d'autre part
\deleted{la}\inserted{l'existence de la} suite \uline{exacte} \unsure{de cofibrations|de cofibration|des cofibrations}, donc une
théorie des cofibres $W$-homotopiques. [NB
Quand la structure est \uline{propre}, i.e. $C\subset\operatorname{Cof}_W$,
alors l'existence de cette structure revient
à dire simplement que ``$W$ admet assez de
$W$-cofibrations'', i.e. que toute $f\in\mathrm{Fl}\,M$
se factorise en $pi$, avec $i\in\operatorname{Cof}_W,p\in W$ ;
et on peut alors prendre $C=\operatorname{Cof}_W$. Dans le
cas où $W$ est stable par \inserted{petites} $\varinjlim$ filtrantes,
$\unsure{C|C_*}$\inserted{$=\operatorname{Cof}_W$}, et $TC=W^{\unsure{\mathrm{univ}|\mathrm{uni}}}$ satisfont \inserted{de plus} l'axiome transfini
$\mathrm{Cof}_4$. \deleted{\unsure{Ces}}] Enfin, K. S. Brown montre
que toute flèche $i\colon X\longrightarrow Y$ dans $C$,
\uline{avec} $X,Y$ \uline{cofibrants} (i.e $\varnothing_M\unsure{-|\longrightarrow}X,\varnothing_M\unsure{\longrightarrow|-}Y$
dans $C$) est dans $\operatorname{Cof}_{W^{\unsure{\natural|\sharp}}}$ (``lemme de
Brown'').

Je m'intéresse ici à la question
d'étendre la structure à cofibrations
à une catégorie de la forme $M(I)$,
et d'utiliser cette extension pour construire
% archive sheet number: 646
